\documentclass[10pt,a4paper,twocolumn]{article}
\usepackage[margin=18mm,columnsep=7mm]{geometry}
\usepackage[T1]{fontenc}
\usepackage[utf8]{inputenc}
\usepackage{lmodern}
\usepackage{microtype}
\usepackage{amsmath,amssymb}
\usepackage{booktabs}
\usepackage{enumitem}
\usepackage{xcolor}
\usepackage{hyperref}
\usepackage{siunitx}
\usepackage{graphicx}
\usepackage{listings}
\usepackage{fancyhdr}

\definecolor{linkblue}{HTML}{075985}
\hypersetup{colorlinks=true,linkcolor=linkblue,citecolor=linkblue,urlcolor=linkblue,pdftitle={Reproducible Terrain-Aware Repeater Field-Strength Mapping with ITU-R P.1812},pdfauthor={Beat W. Meier, HB9MQB}}
\sisetup{detect-all,range-phrase=--,range-units=single}
\DeclareSIUnit{\dBm}{dBm}
\setlist{nosep,leftmargin=*}
\pagestyle{fancy}
\fancyhf{}
\fancyhead[L]{B. W. Meier: Terrain-Aware Repeater Field-Strength Mapping}
\fancyhead[R]{2026}
\fancyfoot[C]{\thepage}
\renewcommand{\headrulewidth}{0.2pt}
\setlength{\parindent}{1em}
\setlength{\parskip}{0pt}

\title{\vspace{-1.2em}\textbf{Reproducible Terrain-Aware Repeater Field-Strength Mapping with ITU-R P.1812}}
\author{Beat W. Meier, HB9MQB}
\date{19 July 2026}

\begin{document}
\twocolumn[
\maketitle
\begin{abstract}
Coverage maps for terrestrial repeaters are often rendered as binary or
single-colour visibility regions. Such products do not preserve a physical
field quantity and cannot distinguish weak from strong predicted service.
This paper specifies an open, reproducible workflow that calculates
path-specific electric field strength in \si{\decibel\micro\volt\per\meter}
using Recommendation ITU-R P.1812-8, terrain profiles over a full
\ang{360}, actual effective radiated power (ERP) where available, and an
explicit P.525 free-space treatment within the short-path region in which a
valid P.1812 profile cannot yet be formed. Polar results are interpolated to an
azimuthal-equidistant raster, reprojected, and published both as numerical and
colour tiles. A cursor service reports field strength and, separately, an
expected nominal S-meter indication under stated ideal receiving assumptions.
The implementation uses Apple M5 performance-core scheduling, shared memory,
and an optional fused MLX/Metal interpolation kernel while retaining a
numerically equivalent NumPy backend. We describe inputs, equations, numerical
contracts, validation layers, uncertainty, and portability. The method improves
information content and reproducibility over binary terrain-visibility maps;
it remains a planning prediction rather than a guarantee of communication.
\end{abstract}
\vspace{0.8em}
\noindent\textbf{Keywords:} radio propagation; field strength; ITU-R P.1812;
terrain profile; repeater; GIS; S meter; Apple silicon; reproducibility
\vspace{1em}
]

\section{Introduction}

Maps are an efficient interface between radio-propagation models and operators,
but the graphical form can obscure the modeled quantity. A terrain horizon or
binary viewshed answers whether two points have geometric visibility under a
particular surface model. It does not calculate received field strength, nor
does it represent diffraction, troposcatter, anomalous propagation, frequency,
ERP, receiver height, radio-climatic zone, or statistical time and location
percentages. A single-colour polygon consequently invites an interpretation
that its boundary separates reliable service from no service.

For terrestrial point-to-area applications from \SI{30}{\mega\hertz} to
\SI{6}{\giga\hertz}, ITU-R P.1812 provides a path-specific prediction method
that combines several propagation mechanisms and returns basic transmission
loss and field strength \cite{iturp1812}. The method builds upon decades of
work on terrestrial propagation and diffraction, including the Bullington and
multiple-knife-edge constructions \cite{bullington1947,deygout1966}.

This work contributes a reproducible computational and publication contract:
(i) physically meaningful \si{\decibel\micro\volt\per\meter} rasters,
(ii) stable numerical web tiles distinct from presentation colours,
(iii) point sampling with an explicitly conditional S-meter estimate,
(iv) provenance manifests, and (v) a native Apple M5 execution strategy with
documented routes to other operating systems and accelerators. The work does
not introduce a new propagation law; its novelty is the auditable integration
of a recognized path model, geospatial sampling, numerical web publication,
and receiver-reference interaction.

\section{Propagation basis}

\subsection{Path-specific field prediction}

Recommendation P.1812-8 predicts terrestrial point-to-area propagation from a
terrain profile and terminal geometry. Its valid frequency interval is
\SIrange{30}{6000}{\mega\hertz}; paths use ground height, representative clutter
height, radio-climatic zones, antenna heights above local ground, polarization,
coordinates, and statistical percentages \cite{iturp1812}. The selected open
implementation is Py1812 at an exact revision rather than an unpinned package
release \cite{py1812}.

For every receiver endpoint, the implementation passes frequency $f$, time
percentage $p$, profile distance $d_i$, terrain $h_i$, representative clutter
$R_i$, climatic zone $z_i$, transmitter and receiver heights $h_{tg}$ and
$h_{rg}$, polarization, terminal coordinates, location percentage $p_L$, and
ERP to P.1812. The model internally accounts for the effective Earth geometry,
diffraction, anomalous/ducting mechanisms and troposcatter. The wrapper must not
add a second Earth-curvature correction.

\subsection{Near-zone treatment}

P.1812 requires an adequate terrain profile; the selected implementation needs
at least five profile points and a path of approximately \SI{0.25}{\kilo\meter}.
Inside this region we use the free-space relation associated with P.525
\cite{iturp525}. Expressed for ERP in watts and distance in kilometres,

\begin{equation}
E_{\mathrm{fs}} = 76.92 + 10\log_{10} P_{\mathrm{ERP}}
                 - 20\log_{10} d
\quad [\mathrm{dB\mu V/m}].
\label{eq:p525}
\end{equation}

This branch is restricted to the short-path domain. It is not a substitute for
P.1812 across the service area. The distinction matters in obstructed terrain,
where free-space extension would systematically overstate distant field.

\subsection{Transmitter and receiver assumptions}

The calculation uses the reported ERP whenever present; \SI{12}{\watt} is an
explicit fallback only for missing values. ERP must not be confused with
transmitter output power before feed-line and antenna effects. The radius is
capped at \SI{100}{\kilo\meter}, receiver height defaults to
\SI{1.5}{\meter} above ground, and the default time/location percentages are
50/50. Vertical polarization is the default but is a station input.

The implementation accepts either a direct non-negative clutter-height raster
or ESA WorldCover classes \cite{zanaga2022worldcover}. The default class mapping
uses representative heights of \SI{15}{\meter} for tree cover,
\SI{3}{\meter} for shrubland, \SI{15}{\meter} for built-up areas,
\SI{1}{\meter} for herbaceous wetland, and \SI{10}{\meter} for mangroves.
The selected mode, source checksums, and mapping are retained in the manifest.
When no clutter input is selected, zero representative clutter remains an
explicit fallback rather than an unrecorded claim of urban accuracy.

\section{Geospatial method}

\subsection{Terrain inputs and vertical meaning}

The engine accepts georeferenced elevation rasters and transforms a bounded
domain into a local azimuthal-equidistant (AEQD) coordinate system centred on
the transmitter. AEQD preserves distance and direction from the centre, which
matches radial path extraction.

Copernicus DEM GLO-30 is a useful worldwide default, but it is a digital surface
model containing vegetation and infrastructure rather than a universal
bare-earth DTM. It has one-arc-second latitude spacing and EGM2008 vertical
heights; reported global aggregate accuracies allow local deviations
\cite{copernicusdem}. Comparative research also demonstrates that DEM errors
depend on terrain and product type \cite{maresova2021}. Therefore, input post
spacing must not be equated with propagation accuracy. Regional DTMs may be
substituted when their datum, license, version, and preprocessing are retained.

\subsection{Polar sampling}

Let $R$ be calculation radius and $\Delta_a$ the requested maximum outer-arc
spacing. With a minimum of 360 rays, the ray count is

\begin{equation}
N_\theta = \max\left(360,
  \left\lceil\frac{2\pi R}{\Delta_a}\right\rceil\right).
\label{eq:rays}
\end{equation}

Azimuths are uniformly distributed in $[0,360)$, avoiding duplication of the
zero-degree ray. Terrain defaults to a \SI{50}{\meter} profile step. Field
endpoints default to \SI{200}{\meter} radial spacing. For every azimuth, the
terrain prefix from transmitter to each endpoint is evaluated by P.1812. This
is computationally expensive but retains the physical endpoint dependence.

Where an externally prepared climatic-zone raster is present, categorical
nearest-neighbour sampling yields sea (1), coastal land (3), or inland (4).
Omitting it explicitly selects inland zone 4 and is inappropriate for rigorous
coastal predictions. For every modeled endpoint the implementation derives the
transmitter- and receiver-terminal distances to the nearest sea transition and
passes them explicitly to P.1812.

\subsection{Rasterization and encoding}

The polar matrix is bilinearly interpolated in azimuth and radius to a default
\SI{60}{\meter} AEQD output grid. This grid spacing is a publication choice,
not a statement that the underlying propagation model resolves every
\SI{60}{\meter} feature. The raster is then reprojected to EPSG:4326.

Numerical PNG tiles use an unsigned 8-bit code $q$. Zero is NoData, and

\begin{equation}
\begin{aligned}
q &= 1 + \operatorname{round}
  \left(\frac{E-E_{\min}}{\Delta E}\right),\\
E_{\min}&=-20,\qquad \Delta E=0.5\ \mathrm{dB},
\end{aligned}
\label{eq:quantize}
\end{equation}

with saturation to $[1,255]$. Presentation tiles apply a fixed absolute scale
from purple through blue, cyan, green, yellow and red. A threshold from 0 to
\SI{50}{\decibel\micro\volt\per\meter}, default 5, changes alpha only. It never
changes stored predictions. Station-independent colours make overlapping
repeaters comparable.

\section{Expected S-meter indication}

The primary output remains $E$ in \si{\decibel\micro\volt\per\meter}. For an
ideal polarization-matched, lossless, \SI{0}{\decibel isotropic} receiving
antenna terminated in a matched receiver, plane-wave power density and
effective aperture yield

\begin{equation}
P_r\,[\mathrm{dBm}] = E\,[\mathrm{dB\mu V/m}]
 +20\log_{10}\lambda\,[\mathrm{m}] -126.755.
\label{eq:fieldpower}
\end{equation}

This is consistent with the free-space transmission framework introduced by
Friis \cite{friis1946}. The nominal IARU Region 1 recommendation assigns S9 to
\SI{-73}{\dBm} below \SI{30}{\mega\hertz} and
\SI{-93}{\dBm} above it, with \SI{6}{\decibel} per S unit below
S9 \cite{iaru_smeter}. Since P.1812 begins at \SI{30}{\mega\hertz}, the latter
reference normally applies here.

Equation~\ref{eq:fieldpower} does not include real antenna gain and direction,
height-dependent local coupling, body loss, polarization mismatch, cable loss,
receiver bandwidth, AGC calibration, or building penetration. The interface
therefore says \emph{expected S-meter indication}, presents the corresponding
field value, and labels the idealized receiver reference.

\section{Implementation and M5 optimization}

P.1812 endpoint calculations are branch-heavy and largely independent by
azimuth. On native Apple silicon, the engine queries the performance-core count
and reserves one or two such cores for the operating system. A process pool
distributes ray chunks. Numeric-library thread counts are limited to one per
worker, avoiding multiplicative oversubscription.

macOS uses the safe \texttt{spawn} start method after GDAL/NumPy initialization.
Naively passing terrain arrays would copy them into every process. The M5 path
places terrain and climatic-zone grids in POSIX shared memory; workers attach
read-only NumPy views. This reduces memory pressure and serialization time.

The polar-to-Cartesian step is regular and massively data-parallel. An optional
MLX backend launches one fused Metal kernel that computes geometry, indices,
four samples, bilinear weights, and the output value. The P.1812 model itself
stays on CPU performance cores. This division matches workload structure and
keeps a NumPy reference backend for numerical comparison. Docker Desktop Linux
guests cannot access the host Metal device, so GPU computation requires a
native macOS process.

\begin{table}[t]
\centering
\caption{Portable acceleration roadmap.}
\label{tab:ports}
\begin{tabular}{p{0.19\columnwidth}p{0.65\columnwidth}}
\toprule
Target & Suggested implementation and evidence \\
\midrule
Linux & NUMA-aware pools, shared memory, SIMD profiling \\
Windows & Spawn-safe mappings and Job Object resource limits \\
NVIDIA & CUDA/CuPy interpolation and colour kernels \\
AMD & ROCm/HIP implementation of the same kernel contract \\
Intel & oneAPI/SYCL on integrated and discrete GPUs \\
Cross-vendor & Vulkan compute; WebGPU for client composition \\
\bottomrule
\end{tabular}
\end{table}

Every accelerator must be compared with the NumPy backend using maximum
absolute and relative error, and must report hardware, OS, runtime, wall time,
peak memory, and energy where available. Moving the whole P.1812 algorithm to a
GPU is not automatically beneficial because divergence and many short profiles
can negate throughput advantages.

\section{API and reproducibility contract}

Each station artifact contains a floating-point GeoTIFF and a JSON manifest.
The manifest records model identifier, pinned implementation, algorithm
version, physical units, coordinates, frequency, ERP, antenna height,
polarization, radius, profile/radial/output spacing, ray count, worker count,
raster backend, field range, and assumptions. Input rasters should additionally
be identified by checksums, vertical datum and license in production workflows.

A stateless read-only HTTP service exposes the manifest, thresholded colour
tiles, immutable numerical tiles, and point samples. The service reads artifact
files rather than exposing an operational database to browsers. This separation
supports TLS termination, caching, rate limiting, atomic artifact replacement,
and independent calculator credentials.

The TypeScript client keeps one imagery layer per station. Filter changes map
to the exact selected station set. Layer toggles and base-map changes do not
invoke camera movement. Pointer sampling is debounced and cancellable; when
several visible stations intersect, the strongest modeled field is shown.

\section{Validation}

Three validation levels must remain distinct.

\paragraph{Analytical and encoding tests.}
P.525 reference values, field/power conversion, S thresholds, tile geometry,
NoData, quantization round trips, and azimuth wraparound are deterministic unit
tests.

\paragraph{Implementation conformance.}
The pinned Py1812 validation profiles should reproduce the recommendation's
reference outputs within their specified tolerance. Synthetic flat, ridge, and
coastal profiles exercise wrapper logic. NumPy and Metal rasterizers must agree
within declared floating-point tolerances. Conformance demonstrates consistent
implementation, not predictive truth at a location.

\paragraph{Empirical validation.}
Calibrated field surveys require known antenna factor, receiver bandwidth,
position, time, transmitter ERP, antenna pattern, cable loss, and preferably
pre-registered routes. Calibration and validation routes should be disjoint.
Report bias, mean absolute error, root-mean-square error, standard deviation and
residual quantiles by distance, terrain class, visibility state and field band.
Uncalibrated consumer S meters are useful qualitative observations but are not
absolute field instruments.

\section{Comparison with binary visibility maps}

Compared with a single-colour terrain-visibility footprint, the specified
method offers six concrete advantages:

\begin{enumerate}
\item the map retains a physical quantity and absolute units;
\item frequency, ERP, heights, polarization, time and location enter the model;
\item diffraction and longer-range P.1812 mechanisms extend beyond geometry;
\item thresholds can change without recalculating or recolouring the physics;
\item numerical tiles support composition and point interrogation; and
\item manifests make assumptions and revisions auditable.
\end{enumerate}

These benefits do not remove uncertainty. A simple visibility map may be faster,
easier to explain, and sufficient for preliminary siting. The appropriate claim
is improved information content and physical traceability, not guaranteed local
reception or universal superiority.

\section{Limitations and future work}

Current limitations include isotropic transmitter patterns, a zero-clutter
fallback when no classified source is supplied, a median-statistics default,
no indoor model, and possible DEM or representative-clutter-height bias.
Coastal predictions require a verified zone raster. Future work should add
antenna-pattern rotation, regional clutter calibration, uncertainty bands,
field-survey tooling, residual dashboards, and sensitivity analysis across
DEMs and statistical percentages.

Software acceleration should expand through tested CUDA, ROCm, SYCL and Vulkan
backends. WebGPU can combine multiple numerical tiles locally and reduce API
load. Such work must preserve the field-value contract and publish numerical
equivalence tests.

\section{Conclusion}

Terrain-aware repeater mapping is most useful when it preserves what was
calculated. The method presented here produces continuous P.1812 field strength
in \si{\decibel\micro\volt\per\meter}, uses an explicit P.525 near zone,
separates numerical values from presentation colours, and makes any S-meter
indication conditional on stated receiving assumptions. Provenance manifests
and layered validation prevent a visually attractive overlay from being
mistaken for measured truth. Apple M5 optimizations reduce runtime without
changing the scientific contract, while the reference NumPy path and documented
porting interfaces support independent implementations.

\section*{Software and data availability}

The source code, web client, API, tests, LaTeX manuscript and citation metadata
are released under the MIT License. Py1812, ITU digital products, elevation data
and map imagery retain their separate licenses and are not redistributed. A
versioned archive should be cited for scientific use.

\section*{Acknowledgements}

The author acknowledges the ITU-R Study Group 3 propagation work, the Py1812
contributors, the IARU technical recommendations, and the providers and
validators of global elevation data.

\bibliographystyle{plainurl}
\bibliography{references}

\end{document}
